\section*{Hoofdstuk \ref{ch:g4:changes-of-state} --- Smelten en koken: veranderingen van toestand}

\begin{solution}{exo:g4:changes-of-state:1}
\omterm{def:g2:ice-water-steam:meltfreeze}{Smelten}; stollen (\omterm{def:g2:ice-water-steam:meltfreeze}{bevriezen}); koken; \omterm{def:g2:ice-water-steam:evapcond}{condenseren}.
\end{solution}

\begin{solution}{exo:g4:changes-of-state:2}
\omterm{def:g4:changes-of-state:melting-point}{Smeltpunt} \qty{0}{\celsius}, \omterm{def:g4:changes-of-state:boiling-point}{kookpunt} \qty{100}{\celsius}. Chocola
\omterm{def:g2:ice-water-steam:meltfreeze}{smelt} rond \qty{34}{\celsius} --- mondtemperatuur.
\end{solution}

\begin{solution}{exo:g4:changes-of-state:3}
\qty{0}{\celsius} --- en vijf minuten later nog steeds
\qty{0}{\celsius}: zolang ijs en water samenwonen, staat de
\omterm{def:g3:temperature-thermometers:temperature}{temperatuur} vastgepind op het \omterm{def:g4:changes-of-state:melting-point}{smeltpunt}.
\end{solution}

\begin{solution}{exo:g4:changes-of-state:4}
Terwijl een stof \omterm{def:g2:ice-water-steam:meltfreeze}{smelt} of kookt, staat haar \omterm{def:g3:temperature-thermometers:temperature}{temperatuur} stil op de
mijlpaal. De warmte gaat helemaal op aan de toestandsverandering zelf
--- \omterm{def:g3:solids-liquids-gases:liquid}{vloeistof} in damp veranderen --- en niet aan klimmen.
\end{solution}

\begin{solution}{exo:g4:changes-of-state:5}
Vast bij \qty{20}{\celsius} (onder $60$); vloeibaar bij
\qty{80}{\celsius} (boven $60$); nog steeds vast bij
\qty{35}{\celsius} --- terwijl een reep chocola in je zak daar al
gesmolten zou zijn.
\end{solution}

\begin{solution}{exo:g4:changes-of-state:6}
Niet heter: beide pannen staan op \qty{100}{\celsius}, het
kookplateau. De extra warmte maakt alleen sneller damp --- verspilde
vlam, dezelfde pasta.
\end{solution}

\begin{solution}{exo:g4:changes-of-state:7}
Aan de twee vlakke stukken --- de plateaus. Overal waar de kromme
horizontaal loopt terwijl het fornuis blijft verwarmen, is een
toestandsverandering aan de gang.
\end{solution}

\begin{solution}{exo:g4:changes-of-state:8}
Goud ($1064 < 1200$) schenk je uit; ijzer ($1538 > 1200$) blijft
vast. De oven komt bij de mijlpaal van ijzer $1538 - 1200 = 338$
graden tekort.
\end{solution}

\begin{solution}{exo:g4:changes-of-state:9}
Smeltend ijs pint het drankje vast op \qty{0}{\celsius}: alle
binnenkomende warmte gaat op aan het \omterm{def:g2:ice-water-steam:meltfreeze}{smelten} van ijs, en niets aan het
opwarmen van het drankje. Pas als de laatste schilfer weg is, kan de
\omterm{def:g3:temperature-thermometers:temperature}{temperatuur} beginnen te klimmen.
\end{solution}

\begin{solution}{exo:g4:changes-of-state:10}
Haar beste gok: de stof is ijs --- bevroren water. De wet dat
mijlpalen kenmerken zijn: het smelt- en \omterm{def:g4:changes-of-state:boiling-point}{kookpunt} van een zuivere
stof veranderen nooit, dus twee getallen die precies met die van water
kloppen zijn zo goed als een handtekening.
\end{solution}

\begin{solution}{exo:g4:changes-of-state:11}
De zekerheid dat water altijd bij \qty{100}{\celsius} kookt. Hoog op
een top kookt het ruim daaronder --- heet, maar geen
\qty{100}{\celsius} --- zodat het ``kokende'' water te koel is om thee
goed te laten trekken. De boosdoener is de ijlere \omterm{def:g2:air-around-us:air}{lucht} die minder op
het water drukt, zoals het hoofdstuk over luchtdruk zal laten zien.
\end{solution}
